% !TEX program = xelatex
\documentclass{英语试卷}
\begin{document}
\examtitle{Z20 名校联盟(浙江省名校新高考研究联盟) 2025 届高三第二次联考}{英语}
\par Z20 名校联盟(浙江省名校新高考研究联盟) 2025 届高三第二次联考\par
\par 英语学科 试题卷\par
\par 命题：黄岩中学 徐彬智、王妮娜\par
\par 磨题：桐乡高级中学 徐敏华 玉环中学 董晓敏 江山中学 毛宇倩 校稿：殷沈华、赵琦\par
\begin{notice}
\par 1. 答卷前，考生务必将自己的姓名、准考证号填写在答题卡上，\par
\par 2. 回答选择题时，选出每小题答案后，用铅笔把答题卡上对应题目的答案标号涂黑。如需改动，用橡皮擦干净后，再选涂其他答案标号。回答非选择题时，将答案写在答题卡上，写在本试卷上无效，\par
\par 3. 考试结束后，将本试卷和答题卡一并交回。\par
\end{notice}
\parttitle{第一部分 听力（共两节，满分 30 分）}
\par 做题时，先将答案标在试卷上，录音内容结束后，你将有两分钟的时间将试卷上的答案转涂到答题卡上。\par
\sectitle{第一节（共5小题：每小题1.5分，满分7.5分）}
\par 听下面 5 段对话。每段对话后有一个小题，从题中所给的 A、B、C 三个选项中选出最佳选项。听完每段对话后，你都有 10 秒钟的时间来回答有关小题和阅读下一小题。每段对话仅读一遍。\par
\question{1. What's the weather like now?}
\begin{choices}(3)
  \item Sunny.
  \item Rainy.
  \item Cloudy.
\end{choices}
\question{2. What is the probable relationship between the speakers?}
\begin{choices}(3)
  \item Co-workers.
  \item Business partners.
  \item Friends.
\end{choices}
\question{3. When will the speakers get there?}
\begin{choices}(3)
  \item At 9:45 a. m.
  \item At 9:30 a. m.
  \item At 9:15 a. m.
\end{choices}
\question{4. Why was the woman late?}
\begin{choices}(3)
  \item She missed the taxi.
  \item The taxi had an accident.
  \item She got injured.
\end{choices}
\question{5. Who will begin the lecture?}
\begin{choices}(3)
  \item Dr. White.
  \item Dr. Hopkins.
  \item Prof. Pears.
\end{choices}
\sectitle{第二节（共 15 小题：每小题 1.5 分，满分 22.5 分）}
\par 听下面 5 段对话或独白。每段对话或独白后有几个小题，从题中所给的 A、B、C 三个选项中选出最佳选项。听每段对话或独白前，你将有时间阅读各个小题，每小题 5 秒钟；听完后，各小题将给出 5 秒钟的作答时间。每段对话或独白读两遍。\par
\par 听第 6 段材料，回答第 6、7 题。\par
\question{6. Where are the speakers?}
\begin{choices}(3)
  \item In a hotel.
  \item In the man's office.
  \item In a post office
\end{choices}
\question{7. What will the man do next?}
\begin{choices}(3)
  \item Fetch an envelope.
  \item Send some papers by fax.
  \item Go downstairs.
\end{choices}
\par 听第 7 段材料，回答第 8 至 9 题。\par
\question{8. What shape of glasses does the man recommend?}
\begin{choices}(3)
  \item Square.
  \item Round.
  \item Heart-shaped.
\end{choices}
\question{9. Which glasses are the least popular?}
\begin{choices}(3)
  \item The red ones.
  \item The gold ones.
  \item The purple ones.
\end{choices}
\par 听第 8 段材料，回答第 10 至 12 题。\par
\question{10. Why does the man go to the woman?}
\begin{choices}(3)
  \item To borrow books.
  \item To submit his work.
  \item To ask for time extension.
\end{choices}
\question{11. Whose book is the man reading now?}
\begin{choices}(3)
  \item Oliver's.
  \item Andrew's.
  \item Billy's.
\end{choices}
\question{12. What will the man probably do?}
\begin{choices}(3)
  \item Design a survey.
  \item Analyze research data.
  \item Interview specialists.
\end{choices}
\par 听第 9 段材料，回答第 13 至 16 题。\par
\question{13. Where is the man from?}
\begin{choices}(3)
  \item Canada.
  \item The UK.
  \item The USA.
\end{choices}
\question{14. What does the man say about bank holidays?}
\begin{choices}(3)
  \item People don't deal with business on these days.
  \item The banks don't open on these days.
  \item The bank services are free on these days
\end{choices}
\question{15. How do people celebrate Shakespeare Day?}
\begin{choices}(2)
  \item Enjoy a related movie. 3. Hold a special party
  \item Record Shakespeare's words.
\end{choices}
\question{16. When did the man last celebrate Shakespeare Day?}
\begin{choices}(3)
  \item Two years ago.
  \item Last year.
  \item Last month.
\end{choices}
\par 听第 10 段材料，回答第 17 至 20 题。\par
\question{17. What was the cause of the disaster?}
\begin{choices}(3)
  \item Too much rainfall.
  \item A stream of mud.
  \item Overflowing rivers.
\end{choices}
\question{18. What were people doing when the disaster hit the area?}
\begin{choices}(3)
  \item Working.
  \item Sleeping.
  \item Traveling.
\end{choices}
\question{19. How many soldiers and police officers came to people's rescue?}
\begin{choices}(3)
  \item Over 100.
  \item Over 200.
  \item Over 1,000.
\end{choices}
\question{20. What might the road conditions lead to?}
\begin{choices}(3)
  \item The difficulty in reaching the trapped.
  \item The possibility of creating anxiety.
  \item The shortage of food and supplies.
\end{choices}
\parttitle{第二部分 阅读理解（共两节，满分 50 分）}
\sectitle{第一节（共15小题：每小题2.5分，满分37.5分）}
\par 阅读下列短文，从每题所给的 A、B、C 和 D 四个选项中，选出最佳答案。\par
\passagelabel{A}
\par Let's get the bad news out of the way first: Dried fruit is a concentrated source of calories and sugars. The keys to getting the most out of it are to keep portions(份额) small and avoid those with added sugars. And choose the fruits that offer the biggest benefits, like the three here.\par
\par 1. Apricots\par
\par PORTION SIZE 5 apricots CALORIES 95\par
\par These orange orbs supply vitamin E, which helps to protect cells against the kind of damage that can trigger cancer and other chronic diseases. You also get 3 grams of fiber and 10 percent of the Daily Value,\par
\par 2. Figs\par
\par PORTION SIZE 5 figs CALORIES 105\par
\par Figs have more fiber than many other dried fruits, and they supply a decent amount of calcium. They may sometimes look like they're sprinkled with sugar, but they're not. The natural fruit sugar in the figs can crystallize and leave a white coating.\par
\par 3. Prunes\par
\par PORTION SIZE 4 prunes CALORIES 92\par
\par Research suggests the mix of antioxidants and other nutrients in prunes may help build up bones, helping to prevent osteoporosis. One serving of prunes has about 20 percent of the Daily Value for vitamin K. The low intakes have been linked to painful blood vessel (血管) problems.\par
\question{21. How many calories will you get if you have a portion of apricots and figs?}
\begin{choices}(2)
  \item 116.
  \item 118.
  \item 124.
  \item 200.
\end{choices}
\question{22. What is a benefit of prunes?}
\begin{choices}(2)
  \item Relieve pain.
  \item Improve bone strength.
  \item Prevent cancer.
  \item Promote calcium absorption.
\end{choices}
\question{23. What advice is offered on dried fruit?}
\begin{choices}(2)
  \item Take small servings.
  \item Add some sugar.
  \item Mix with other nutrients.
  \item Stay away from it.
\end{choices}
\passagelabel{B}
\par Walking home from school, I am enjoying my day. Each day I stop to hear music coming from the ivy(常春藤). Sparrow(麻雀) song, a loud glorious sound. The ivy is pulsing below a leafless tree, decorating around the lower trunk like a celebration. Nervous flittings, in and out, pecking at the branches. This ivy is a home for them, harboring an entire flock. But this scene is no longer ubiquitous. Numbers have declined by almost 70\% in the UK since 1970. Sparrows make their homes near human places and finding a tree like this is a blessing. Happily, the ivy is still full of chatter and twitter.\par
\par In Greek mythology, sparrows are sacred and often associated with the goddess Aphrodite, symbolizing true love and spiritual connection. I wonder how many people look at sparrows and feel that depth of connection. All birds live brightly in our imagination, connecting us to the natural world, opening up all kinds of creativity.\par
\par As I stand there, fluffed-up sparrows in conversation among themselves, there is a spark. Noticing nature is the start of it all. Slow down to listen, to watch. Take the time, despite mountains of homework.\par
\par Take the time to observe nature, to immerse oneself in its patterns, structures, happenings and rhythms. It is how mathematicians and scientists are nurtured. Alan Turing studied the patterns in nature: the arrangement of petals on the flower, spots on a leopard and stripes on a zebra. He was looking for a mathematical formula(公式) for the development of cells in living things. Nature sparks creativity. All we have to do is start with the question, Why? The way my mind whirrs and whirls in nature, or even when "daydreaming" is way more productive than the work I do in school.\par
\par I consolidate myself by thinking and thinking while intensely watching the movements of sparrows. Who knows where watching sparrows will lead?\par
\question{24. What do I do on the way home from school?}
\begin{choices}(2)
  \item Listen to pop music.
  \item Decorate the tree.
  \item Clean up the trunk.
  \item Observe the sparrows.
\end{choices}
\question{25. What does the underlined word "ubiquitous" in paragraph 1 probably mean?}
\begin{choices}(2)
  \item Unusual.
  \item Delightful.
  \item Common.
  \item Horrible.
\end{choices}
\question{26. What is the author's purpose in mentioning Greek mythology?}
\begin{choices}(2)
  \item To show the imagination of humans.
  \item To stress the significance of sparrows.
  \item To introduce the power of Aphrodite.
  \item To demonstrate the diversity of species.
\end{choices}
\question{27. What can we learn from paragraph 4 of the text?}
\begin{choices}(2)
  \item Nature facilitates cell development.
  \item Nature follows rigid patterns.
  \item Nature is a source of inspiration.
  \item Nature offers companionship.
\end{choices}
\passagelabel{C}
\par After Taylor Swift, the "Mona Lisa" is probably the most recognizable female face in the world. Every day around 20,000 people stare at the Leonardo da Vinci's painting in the Louvre. Yet it became famous not because of a tempting semi-smile. Until a worker stole the masterpiece in 1911, it was still mostly unknown; viewers flooded in to see what a French newspaper called "an enormous, horrific, gaping void(空白)”.\par
\par Classics of art, literature and music are supposed to carry some mysterious appeal that endures across the ages. But as Rochelle Gurstein, a historian, writes in a new book, the "timeless classic" is anything but. "What I believed was written in stone was actually written in water", she argues. Classics come and go.\par
\par Take other celebrated works of art. Until the end of the 19th century Michelangelo's "David" had only his fig leaf to protect him from the weather in Piazza della Signoria in Florence. He attracted more pigeons than visitors; visitors preferred the artist's sculpture of Moses in Rome. Before the famous "Venus de Milo" was discovered in 1820, viewers were crazy about "Venus de Medici", a delicate sculpture most people alive today have never heard of.\par
\par Much as in art, the classics of music and literature have not always been set in stone either. Today, Johann Scbastjan Bach is often regarded as the master of classical music. But after his death in 1750, Bach's compositions fell out of fashion until Felix Mendelssohn, a German composer, rediscovered them in the 1820s. Herman Melville's masterpiece "Moby Dick" was largely forgotten until critics such as Raymond Weaver, a professor at Columbia University rescued him from the depth of a nobody around the centenary of his birth in 1919.\par
\par "The status of classics" depends on critics and shows the power of passionate writers to change the way people look at art. Tastes evolve, and they should. But culture lives on only as long as people continue to argue over what counts as a classic today. The true enemy of the classic is not relativism or awakening. It is silence.\par
\question{28. What made the painting "Mona Lisa" famous?}
\begin{choices}(2)
  \item The theft.
  \item The semi-smile.
  \item The report.
  \item The painter.
\end{choices}
\question{29. What does Rochelle think of classics?}
\begin{choices}(2)
  \item Enduring.
  \item Fading.
  \item Changeable.
  \item Admirable.
\end{choices}
\question{30. How does the author develop paragraph 4?}
\begin{choices}(2)
  \item By making a comparison.
  \item By offering examples.
  \item By giving an explanation.
  \item By giving a description.
\end{choices}
\question{31. What might ultimately discourage the development of classics?}
\begin{choices}(2)
  \item The ignorance of culture.
  \item The course of time.
  \item The evolution of tastes.
  \item The lack of criticism
\end{choices}
\passagelabel{D}
\par The brain is protected by shock-absorbing fluid and, outside that, the skull. In a concussion(脑震荡), the brain bounces around in the skull, accelerating, decelerating or revolving. This creates a series of impacts. The neurons(神经元) in the brain are disturbed, and between those neurons, the axons - thin fibers that transmit electrical impulses – stretch or break. Inside the brain, it's like an earthquake that has happened. Everything is still standing, but there\par
\par are cracks in the roads and in building foundations.\par
\par This extremely small damage has big consequences. But decades ago, the public had a "walk it off" attitude towards concussions. That's changed, thanks in part to an increased awareness about the long-term effects. The 2015 film "Concussion" tells the true story of Dr Bennet Omalu, a forensic pathologist studying chronic traumatic encephalopathy (CTE), a brain disorder caused by repeated head injuries.\par
\par Until recently, the dark-room treatment was the standard prescribed by most doctors, involving patients resting in darkness and without any mental stimulation, until their symptoms eased off. This suggestion came out of studies in the 1990s and early 2000s that found when athletes with concussions continued with high levels of activity, they recovered more slowly than those who rested.\par
\par But people who spend more time in a dark room - sometimes called "cocooning" - are also more likely to experience anxiety, depression and sleep problems. Their bodies get weaker due to underuse, as well.\par
\par "Science and research have evolved over the past decade, showing that when you stop doing activity altogether, that actually affects your recovery process," says Babul, a sports-injury specialist. The active-rest approach to concussion recovery promotes a slow increase in activity, starting with a day or two of light daily movement, and gradually adding activity, building to a full return to normal, unrestricted activity.\par
\par "People who follow the right guidance and management tend to recover uneventfully, and those who don't and continue to do their activities continue to tax their brain and are likely to have long-term, persisting symptoms," Babul says. "The key is to recognize it immediately and know what to do."\par
\question{32. What is the first paragraph mainly about?}
\begin{choices}(2)
  \item The workings of concussion.
  \item The causes of concussion.
  \item The symptoms of concussion.
  \item The effects of concussion.
\end{choices}
\question{33. What leads to people's weaker bodies in the dark-room treatment?}
\begin{choices}(2)
  \item The dark environment.
  \item Their physical pain.
  \item Their insufficient activity.
  \item Their mental disorder.
\end{choices}
\question{34. Which practice is advised to help with concussion, according to the passage?}
\begin{choices}(2)
  \item Setting a progressive fitness plan first.
  \item Doing heavy housework daily.
  \item Playing competitive sports regularly.
  \item Lying in bed for the whole day.
\end{choices}
\question{35. What is the best title of the text?}
\begin{choices}(2)
  \item From Inaction to Action: Concussion Shocks
  \item From Darkness to Light: Concussion Recovery
  \item From Outside to Inside: Concussion Treatment
  \item From Nothing to Something: Concussion Study
\end{choices}
\sectitle{第二节（共5小题；每小题2.5分，满分12.5分）}
\par 阅读下面短文，从短文后的选项中选出能填入空白处的最佳选项。选项中有两项为多余选项，\par
\par In the freezing cold winter, air conditioners and heaters work overtime. The satisfaction of feeling a warm breeze against your skin is an amazing experience. \blanknum{36}\par
\par What are to blame are often the inefficient home appliances you have. Luckily, there are smart home appliances that can alleviate this burden. These appliances offer many benefits, all while keeping your electric bill manageable.\par
\par \blanknum{37} Smart lighting systems can detect activity in an area and adjust illumination accordingly, providing energy efficiency. In terms of security, smart door locks offer peace of mind by automatically securing your home when you leave.\par
\par And smart appliances can do more. From floor-cleaning robots to smart laundry systems, these appliances are equipped with advanced sensors. They make household chores effortless, leaving you more time for recreational activities. \blanknum{38} You can even set a romantic mood in the living room by adjusting the smart lights and playing\par
\par a romantic comedy movie on the smart TV.\par
\par Now, if you're considering switching to smart home appliances, it's advisable to approach this gradually. \blanknum{39} Smart lights, switches, plugs, smart hubs, and of course, Wi-Fi are wise choices. These small items play a crucial role in your smart home setup, especially the smart hubs and the Wi-Fi.\par
\par Switching from traditional and inefficient home appliances to smart ones is more than just a modern convenience. \blanknum{40} You're not just simplifying your daily routines but actively contributing to environmental conservation efforts.\par
\begin{choices}(1)
  \item But it needs a greener solution.
  \item But the following electric bill can be a shock.
  \item It's a significant step towards a sustainable future.
  \item Smart TVs provide many entertainment options for you to explore.
  \item Moreover, smart appliances extend their features beyond the chores.
  \item Advanced appliances can effectively manage every aspect of your household.
  \item As large appliances can be significant investments, it is easier to start with smaller items.
\end{choices}
\parttitle{第三部分 语言知识运用（共两节，满分 30 分）}
\sectitle{第一节（共15小题；每小题1分，满分15分）}
\par 阅读下面短文，从短文后各题所给的 A、B、C 和 D 四个选项中，选出可以填入空白处的最佳选项。\par
\par Six years as a space reporter taught me that chaos rules everything. Nothing is completely predictable or controlled. I have watched enough go wrong to know that no view of a celestial(天空的) event is ever truly promised to us \blanknum{41} on the earth.\par
\par So when my mother and I decided to drive out to Erie, Pennsylvania, to see the April 8 total solar eclipse(日全食) this year, I knew this would be a trip of two \blanknum{42} attitudes. I'd be \blanknum{43} about the weather and convinced we'd be victims of the randomness that governs the world while my mom would have strong faith that order would \blanknum{44} and the skies would let us glimpse an eclipse like this for the first time in our lives. I told her not to \blanknum{45} the universe for this one; she told me she wouldn't depend on anything else.\par
\par I was once again humbled into a \blanknum{46} I've learned time and time again: Mama knows best. Erie's forecast that morning was looking \blanknum{47}, but by the time first contact between the sun and the moon began a little after 2 p. m., the clouds over the city's bay began to clear. The pale yellow sun under the eclipse lenses rapidly crested(达到顶峰), concentrating into a \blanknum{48} orange glow.\par
\par The eclipse \blanknum{49} at 3:16 p. m. A thin white glow pierced out from the edge of a clean black circle. The sunset-like \blanknum{50} glowed and bloomed in the distance. I could spot solar prominences(H 珩) sticking out from the sides of the sun \blanknum{51} slightly bright red and pink. Jupiter and Venus also played roles in this performance. It felt like bearing \blanknum{52} to something close to a miracle.\par
\par Four minutes later, the eclipse \blanknum{53}. The sun brightened again. And the clouds \blanknum{54} violently, swallowing up the moon and the sun and the sky in gray. But for four incredible minutes, the universe seems to have \blanknum{55} a promise to my mother.\par
\cz{41}{analysts}{observers}{investigators}{astronauts}
\cz{42}{consistent}{conflicting}{amusing}{agreeable}
\cz{43}{optimistic}{confused}{pessimistic}{curious}
\cz{44}{win}{lose}{delay}{fail}
\cz{45}{feed on}{count on}{take on}{settle on}
\cz{46}{plan}{skill}{lesson}{practice}
\cz{47}{usual}{average}{beautiful}{terrible}
\cz{48}{fierce}{weak}{hazy}{dim}
\cz{49}{disappeared}{exploded}{died}{struck}
\cz{50}{shapes}{sizes}{colors}{scenes}
\cz{51}{in}{on}{with}{for}
\begin{choices}(3)
  \item tolerance
  \item relation C: witness
  \item reference
\end{choices}
\cz{53}{continued}{progressed}{ticked}{ended}
\cz{54}{quit}{returned}{left}{faded}
\cz{55}{made good on}{lost focus on}{took account of}{kept track of}
\sectitle{第二节（共10小题；每小题1.5分，满分15分）}
\par 阅读下面短文，在空白处填入 1 个适当的单词或括号内单词的正确形式，\par
\par Su embroidery(刺绣), acknowledged as the finest Chinese hand needlework, is an ancient craft that originated from Suzhou, a city with a long history of silk making. The garden city of Suzhou was \blanknum{56} economic and cultural center in ancient China, where the women of noble families in particular had a taste \blanknum{57} art and their needlework style was influenced by their aesthetic tastes.\par
\par "\blanknum{58} (feature) its rich colors, complex techniques and elegant designs, the craft \blanknum{59} (reflect) the beautiful natural environment and the leisurely lifestyle of Suzhou since centuries ago," said Zheng Lihong, an art professor at Soochow University in Jiangsu province. "It is an art form created by women living in a city \blanknum{60} prides itself on an artistic lifestyle." Zheng added.\par
\par The embroidery reached \blanknum{61} (it) golden age during the Song Dynasty, integrating elements of the ancient paintings popular at the time. The craft sbecame a favorite of Chinese \blanknum{62} (royal) in the court of the Qing Dynasty. It was during this period that the term suxiu, Su embroidery, \blanknum{63} (coin). In 1954 the Suzhou Silk Embroidery Research Institute was established, which made it possible for many innovations in silk embroidery techniques \blanknum{64} (follow) soon.\par
\par Zou Yingzi, an inheritor of Su embroidery, can \blanknum{65} (flexible) apply various stitching techniques to different parts of her creations, whose works are now widely collected by museums in China.\par
\parttitle{第四部分 写作（共两节，满分 40 分）}
\sectitle{第一节（满分15分）}
\par 假设你是李华，你校下周日有两个志愿者活动同时举行，一个在火车站，另一个在市文化博物馆。你的交换生朋友 Peter 不知道选择哪一个参加，写信向你寻求建议。请你写一封回信，内容包括：\par
\par 1. 给出建议；\par
\par 2. 说明理由。\par
\par 注意:\par
\par (1) 写作词数应为 80 左右;\par
\par (2) 请按如下格式在答题纸的相应位置作答，\par
\par Dear Peter, Yours, Li Hua\par
\sectitle{第二节（满分25分）}
\par 阅读下面材料，根据其内容和所给段落开头语续写两段，使之构成一篇完整的短文。\par
\par Every October, I beg my parents to take us to a pumpkin patch(园地) and every October it doesn't work out. That's why I couldn't believe my ears when my dad came home with the world's best news.\par
\par "Guess what!" he said. "I just saw a poster for the jack-o’-lantern(南瓜灯) design party in the Pumpkin Patch, and it’s not too far from here."\par
\par "No way!" I said. "Can we go?"\par
\par "Pumpkin! Pumpkin! Party! Party!" chanted my little brother, Tex.\par
\par "Let's go! But I've heard a rumor that something terrible happened there." said my dad, worried.\par
\par Within minutes, Tex and I already jumped into the car and buckled up, ready to go. "I'm going to pick the giantest pumpkin," said Tex.\par
\par "Giantest isn't an actual word," I said.\par
\par "I'm going to pick the bestest pumpkin!" shouted Tex.\par
\par "I'm pretty sure bestest isn't an actual word either." I giggled. "I'm going to pick whichever pumpkin will make the scariest jack-o'-lantern."\par
\par "All right, boys! Too much excitement! What if something really goes wrong?" my dad cut in.\par
\par No response. My dad continued, "If plan A fails, we can get creative and come up with Plan B!"\par
\par Still no response.\par
\par Tex and I continued our excitement, giggling and chatting.\par
\par The drive was extra cool because we saw all kinds of Halloween decorations along the way.\par
\par "Look!" Tex said, pointing to twinkling orange lights and a flashing Pumpkin Patch sign.\par
\par "Yeah! Yes! It wasn't too far a drive," I said. "Watch out, pumpkins! Here we come!"\par
\par "I pick that giant one!" Tex said, running over to pick it up. But just before he did, he said, "Ew! There's a hole in it!" Tex quickly raced to another perfectly perfect pumpkin. But when he got closer, he said, "No fair! Someone Coming into our sight were pumpkins of various sizes that lay in the patch. took a bite."\par
\par 注意:\par
\par (1) 续写词数应为 150 左右;\par
\par (2) 请按如下格式在答题卡的相应位置作答。\par
\writingpara{Paragraph 1: That's when I noticed the Pumpkin Patch Closed sign.}
\writinglines{8}
\writingpara{Paragraph 2: Finally came our own jack-o'-lantern design party the next day.}
\writinglines{8}
\end{document}
